\section{Phương pháp kiểm thử và đánh giá}

Phần này chỉ xác định chiến lược kiểm thử dùng cho mini project. Các kết quả kiểm thử cụ thể được trình bày ở Chương 3, còn logic chi tiết của Rule Evaluation Engine, API, cache, RBAC và audit log đã được mô tả trong các phần thiết kế tương ứng.

\subsection{Phạm vi kiểm thử}

Hệ thống được kiểm thử theo bốn lớp chính:

\begin{itemize}
    \item \textbf{Domain và Rule Evaluation Engine}: kiểm tra precedence, targeting rule, percentage rollout, kill switch, default off và fallback.
    \item \textbf{Backend API}: kiểm tra contract của Management API và Evaluation API, bao gồm validation, lỗi không tìm thấy, xung đột dữ liệu và response shape.
    \item \textbf{Control-plane side effects}: kiểm tra audit log, cache invalidation, thống kê evaluation và phân quyền demo.
    \item \textbf{Workflow trình diễn}: kiểm tra Dashboard, Demo Application, JavaScript SDK và Docker Compose trên các kịch bản có thể trình bày trực tiếp.
\end{itemize}

Cách chia này giúp tránh kiểm thử trùng lặp cùng một hành vi ở nhiều tầng. Unit test tập trung vào quyết định thuần của engine; integration test xác nhận API, persistence và side effect; kịch bản demo xác nhận trải nghiệm end-to-end.

\subsection{Tiêu chí đánh giá Rule Evaluation Engine}

Rule Evaluation Engine được đánh giá theo ba tiêu chí:

\begin{itemize}
    \item \textbf{Đúng precedence}: trạng thái archive, disabled, group kill switch, flag kill switch và global on phải có thứ tự ưu tiên cao hơn các targeting rule.
    \item \textbf{Xác định}: cùng một snapshot cấu hình và cùng một evaluation context phải luôn tạo ra cùng một kết quả, đặc biệt đối với percentage rollout.
    \item \textbf{Fail-closed}: khi flag không tồn tại, context không đủ hoặc engine gặp lỗi, hệ thống không được tự động bật tính năng ngoài ý muốn.
\end{itemize}

Các test không lặp lại mô tả thuật toán; chúng chỉ tạo dữ liệu đầu vào, gọi engine hoặc Evaluation API và so sánh với kết quả kỳ vọng gồm \texttt{enabled}, \texttt{variant}, \texttt{reason} và metadata rule nếu có.

\subsection{Tiêu chí đánh giá API và dữ liệu}

Đối với API, trọng tâm kiểm thử là contract và tính nhất quán của dữ liệu. Các mutation phải validate input, ghi dữ liệu đúng ràng buộc, tạo audit log trong cùng transaction và làm mất hiệu lực cache liên quan. Các endpoint đọc phải trả về cấu trúc response ổn định, hỗ trợ phân trang khi cần và không để lộ dữ liệu nội bộ.

Đối với Evaluation API, response luôn phải có dạng evaluation-shaped kể cả khi tài nguyên không tồn tại hoặc xảy ra lỗi. Điều này giúp SDK và Demo Application xử lý fallback bằng cùng một contract thay vì phải phân nhánh theo nhiều kiểu lỗi HTTP khác nhau.

\subsection{Tiêu chí đánh giá Dashboard và Demo Application}

Dashboard được đánh giá theo khả năng hỗ trợ thao tác quản trị an toàn: hiển thị đúng trạng thái lifecycle/runtime, vô hiệu hóa thao tác không đủ quyền, phản ánh lỗi validation và cho phép theo dõi audit log/statistics. Demo Application được đánh giá theo khả năng chứng minh việc thay đổi cấu hình Feature Flag có thể làm thay đổi hành vi ứng dụng mà không cần deploy lại.

Các kịch bản demo được chọn để bao phủ những giá trị quan trọng nhất của sản phẩm: bật/tắt toàn cục, targeting theo user hoặc role, rollout theo phần trăm, flag kill switch, group kill switch, audit log, statistics và fallback của SDK.

\subsection{Tiêu chí triển khai và trình diễn}

Quy trình triển khai được đánh giá bằng khả năng khởi tạo hệ thống trong môi trường sạch. Docker Compose phải khởi chạy được PostgreSQL, Backend API, Admin Dashboard và Demo Application; migration và seed data phải có thể chạy lại theo cách idempotent; các service phải giao tiếp đúng qua cấu hình URL và biến môi trường.

Tiêu chí cuối cùng là khả năng trình bày: sau khi khởi chạy, người xem có thể đi từ Dashboard sang Demo Application, thay đổi cấu hình flag và quan sát kết quả runtime cùng audit/statistics phát sinh.
